\chapter{Einleitung}

\section{Hintergrund und Motivation}

Die Digitalisierung ist heute allgegenwärtig und stellt einen unverzichtbaren Bestandteil 
der Kommunikation zwischen Unternehmen und Kunden dar. Insbesondere im Bereich der 
Humanmedizin wird deutlich, wie digitale Technologien im Kontext der Kommunikation einen 
Vorteil bieten. Unternehmen wie Doctolib haben bereits früh das Ziel gesetzt, die 
Terminverwaltung zwischen Arztpraxen und Patienten durch den Einsatz digitaler 
Terminverwaltungssysteme effizienter und nutzerfreundlicher zu gestalten. Dieses Ziel haben 
sie weitgehend erreicht: Die digitale Terminbuchung ist heute nicht mehr wegzudenken und 
wird von vielen Patienten als entscheidender Faktor für die Inanspruchnahme von ärztlichen 
Leistungen wahrgenommen.

Vor diesem Hintergrund drängt sich die Frage auf, wie die Situation im 
veterinärmedizinischen Bereich aussieht. In Deutschland gibt es rund 10.500 
Tierarztpraxen, von denen über die Hälfte Kleintierpraxen sind, die sich auf Haustiere 
spezialisiert haben (Bundestierärztekammer, 2024). Viele Tierarztpraxen nutzen bereits 
Praxisverwaltungssoftware, um interne Aufgaben wie Terminplanung, Abrechnung und 
Dokumentation zu organisieren. Allerdings existiert bisher kein etabliertes Unternehmen, 
das den Markt für die Terminverwaltung zwischen Tierbesitzern und Tierarztpraxen ähnlich 
beherrscht wie etwa Doctolib in der Humanmedizin.

Zudem gibt es bisher keine repräsentativen Studien, die ermitteln, wie viele 
Tierarztpraxen bereits eine Terminverwaltungssoftware mit Online-Terminbuchung für 
Tierbesitzer in ihren Alltag integriert haben. Dennoch lassen Marktberichte darauf 
schließen, dass sich die Digitalisierung der Veterinärmedizin der der Humanmedizin 
zunehmend annähert: Der globale Markt für Veterinärpraxis-Managementsoftware wird für 
2023/24 auf etwa 1,8 bis 2,0~Milliarden US-Dollar geschätzt und soll bis etwa 2032 auf 
rund 3,2~Milliarden US-Dollar wachsen, mit einer jährlichen Wachstumsrate von etwa 6--7~\%. 
Auf dieser Grundlage lässt sich plausibel argumentieren, dass die Digitalisierung der 
Veterinärbranche sich in jeglichen digitalen Bereichen weiterentwickeln wird -- damit ist 
auch die Terminverwaltung zwischen Tierarztpraxen und Tierbesitzern gemeint.

\section{Zielsetzung}

Vor diesem Hintergrund zielt die vorliegende Bachelorarbeit darauf ab, eine 
Terminverwaltungssoftware zu konzipieren, umzusetzen und zu evaluieren, um 
Tierbesitzerinnen die Terminbuchung zu erleichtern. Die Software soll ermöglichen, 
Haustiere anzulegen, Tierpraxen in der Nähe zu suchen und Termine jederzeit online zu 
buchen. Der zentrale Mehrwert liegt in der flexiblen, ortsunabhängigen Online-Terminvergabe, 
die die Telefonlast reduziert und eine einfache Verwaltung für Nutzende bietet. Im 
Vergleich zu bestehenden Lösungen hebt sich \textit{Petappoint} durch seinen Fokus auf 
die Terminverwaltung ab. Onlineberatungsgespräche oder das Tracken der Tiergesundheit 
sind nicht Teil der Software.

Die Arbeit grenzt sich bewusst ab: Eine Tierpraxenansicht mit Funktionen für das Personal 
wird nicht umgesetzt. Der Schwerpunkt liegt ausschließlich auf der Konzeption, Umsetzung und 
Evaluation der Software aus Sicht der Tierbesitzer.
